\centerline{\bf \Huge Стратегия выигрыша}
\title{Стратегия выигрыша}

\begin{flushright}
    \scriptsize Наблюдай, анализируй,\\ предугадывай действия противника.\\
    Обещанный Неверленд
\end{flushright}

\problem{1.1}Имеется три кучки камней: в первой – 10, во второй – 15, в третьей – 20. За ход разрешается разбить любую кучку на две меньшие. Проигрывает тот, кто не сможет сделать ход. Кто выиграет?

\problem{1.2}Числа от 1 до 20 выписаны в строчку. Игроки по очереди расставляют между ними плюсы и минусы. После того, как все места заполнены, подсчитывается результат. Если он чётен, то выигрывает первый игрок, если нечётен, то второй. Кто выиграет?

\problem{1.3} Двое по очереди ставят ладей на шахматную доску так, чтобы ладьи не били друг друга. Проигрывает тот, кто не может сделать ход. Кто выиграет?

\problem{2.1} На столе лежат две стопки монет: в одной из них 30 монет, а в другой - 20. За ход разрешается взять любое количество монет из одной стопки. Проигрывает тот, кто не сможет сделать ход. Кто из игроков выигрывает при правильной игре?

\problem{2.2}На доске написано число 1. Два игрока по очереди прибавляют любое число от 1 до 5 к числу на доске и записывают вместо него сумму. Выигрывает игрок, который первый запишет на доске число тридцать. Укажите выигрышную стратегию для второго игрока.

\problem{2.3} У ромашки 12 лепестков. За ход разрешается сорвать либо один лепесток, либо два рядом растущих лепестка. Проигрывает игрок, который не сможет сделать ход. Как действовать второму игроку, чтобы выиграть независимо от ходов первого игрока?

\problem{3}Имеется две кучки бэбифоксов - по 2023 в каждой. За ход разрешается взять любое количество бэбифоксов, но только из одной кучки. Проигрывает тот, кому нечего брать.

\problem{4}На столе у Начальника лежит 37 чокопаев. Разрешается по очереди брать не более 5 чокопаев. Выигрывает тот, кто возьмет последнюю. Кто выигрывает при правильной игре?

\newpage

\problem{5}Шахматный король стоит в левом нижнем углу шахматной доски. Участвуют два игрока, которые ходят по очереди. За один ход его можно передвинуть на одно поле вправо, на одно поле вверх или на одно поле по диагонали "вправо-вверх". Выигрывает игрок, который поставит короля в правый верхний угол доски. Кто из игроков выигрывает при правильной игре?

\problem{6}В каждой клетке доски 11 × 11 стоит шашка. За ход разрешается снять с доски любое количество подряд идущих шашек либо из одного вертикального, либо из одного горизонтального ряда. Выигрывает снявший последнюю шашку.  

\problem{7}На доске размером 8×8 двое по очереди закрашивают клетки так, чтобы не появлялось закрашенных уголков из трёх клеток. Проигрывает тот, кто не может сделать ход. Кто выигрывает при правильной игре?

\problem{8} В ряд лежат 25 чупачупсов. За ход разрешается брать один или два рядом лежащих чупачупса. Проигрывает тот, кому нечего брать.

\problem{9} У Начальника есть прямоугольный стол. Два ученика начинают по очереди класть на него по одному рублю так, чтобы эти монеты не перекрывали друг друга. Кто не может сделать ход - проигрывает. Кто выиграет при правильной игре?